

\import{chapters/4. Thiết kế hệ thống E-Commerce tích hợp Pipeline/}{4.1 Đề xuất kiến trúc tổng thể.tex}
\import{chapters/4. Thiết kế hệ thống E-Commerce tích hợp Pipeline/}{4.2 Các luồng thực thi chính.tex}
\import{chapters/4. Thiết kế hệ thống E-Commerce tích hợp Pipeline/}{4.3 Xây dựng và tổ chức dữ liệu.tex}
% \import{chapters/4. Thiết kế hệ thống E-Commerce tích hợp Pipeline/}{4.4 Thiết kế giao diện người dùng.tex}

\section{Kết chương}

Tóm lại, việc mô hình hóa Minesweeper thành một đồ thị vô hướng $G = (V, E)$ và biểu diễn thuật toán DFS bằng cấu trúc Stack không chỉ là lý thuyết, mà là nền tảng để đảm bảo hệ thống chạy ổn định với độ phức tạp $O(N)$ cả về thời gian lẫn bộ nhớ.

Những ràng buộc toán học này đặt ra yêu cầu rõ ràng cho việc hiện thực hóa. Chương tiếp theo sẽ trình bày cách các ý tưởng này được chuyển thành cấu trúc dữ liệu cụ thể và cách thuật toán vận hành.

